% CONCLUSIONES Y TRABAJOS FUTUROS

\section{Conclusiones y trabajos futuros}

La caracterización del campo permitió identificar una regularidad clara en la literatura reciente: la calidad del aire se estudia con combinaciones heterogéneas de variables, arquitecturas y propósitos de uso, pero persisten tensiones en trazabilidad, comparabilidad y explicitud del razonamiento cuando se pasa de modelos aislados a módulos operativos. Esa lectura justificó concentrar la propuesta en los contaminantes criterio del AQI, separar temperatura y humedad como contexto, y adoptar datos abiertos como fuente principal de observación. En ese sentido, la revisión sistemática aportó antecedentes comparables y delimitó criterios de diseño verificables y coherentes con el tipo de problema abordado.

Sobre esa base, el diseño arquitectónico quedó resuelto en una estructura funcional de cinco capas: consolidación normativa, variables auxiliares, inferencia difusa principal, capa contextual basada en reglas y alertamiento. Esta partición permitió traducir el problema en un flujo legible, con entradas, derivaciones y salidas diferenciadas, y dejó explícita la separación entre el cálculo normativo del AQI, la interpretación difusa del riesgo y el ajuste contextual posterior. El resultado es una arquitectura auditable, preparada para mostrar cómo se forma el estado base, cómo se activan las reglas y cómo se compone la salida final sin confundir fuentes de dato, lógica normativa e interpretación experta.

La implementación confirmó que esa arquitectura podía materializarse en un artefacto ejecutable. El prototipo integra adquisición de datos abiertos, cálculo AQI con \textit{breakpoints} \textit{EPA/AQS}, una base principal de 54 reglas difusas, una base contextual crisp de 9 reglas, control de cobertura, API HTTP, ejecución por CLI y Docker, e interfaz web con vistas de \textit{dashboard}, trazabilidad, explicabilidad y evaluación. Esta integración deja como resultado un módulo reproducible y técnicamente coherente, capaz de exponer el valor final del episodio y la lógica que lo produce. La implementación todavía conserva límites operativos en persistencia robusta del histórico local, administración y endurecimiento de pruebas de interfaz, pero ya supera la etapa de prueba conceptual aislada.

La validación realizada en este proyecto muestra que el comportamiento del módulo es consistente con el alcance fijado para la propuesta. Las corridas controladas y las consultas reales permitieron verificar cálculo normativo, activación de reglas, uso de variables auxiliares, funcionamiento de la capa contextual y exposición trazable de la salida. La evidencia reunida sostiene la utilidad técnica del artefacto para evaluación actual del riesgo y generación de alertas sobre datos abiertos. A la vez, esta validación sigue siendo acotada frente a una comparación extensa automatizada entre múltiples estaciones, ventanas temporales y contextos urbanos. Esa restricción define el alcance real de los resultados: existe validación funcional suficiente del módulo, pero no una validación externa amplia comparable a estudios predictivos o de calibración multisitio.

La propuesta cumple el propósito general de desarrollar un módulo de monitoreo de calidad del aire orientado a apoyar la evaluación del riesgo y la generación de alertas mediante inferencia difusa y datos abiertos en entornos urbanos. La contribución principal de esta investigación es la construcción justificada de un artefacto explicable, trazable y ejecutable que integra base normativa, razonamiento difuso y evidencia visual operativa. El alcance presentado corresponde a un módulo funcional de evaluación actual del riesgo; no describe una plataforma productiva cerrada ni un modelo predictivo de largo horizonte. Ese cierre es consistente con el problema formulado, con los objetivos definidos y con la evidencia presentada a lo largo de la memoria.

Como trabajo futuro inmediato, resulta pertinente ampliar la validación comparativa creando funcionalidades para evaluar de manera automática múltiples estaciones, más horizontes temporales y episodios de mayor diversidad, además de reforzar la cobertura de pruebas API, frontend y E2E. También conviene sustituir el histórico local por una persistencia más robusta y, en una etapa posterior, estudiar una variante contextual difusa que permita contrastar su comportamiento con la capa crisp actual. En paralelo, una línea posterior de investigación puede aprovechar el motor de inferencia construido en esta propuesta para generar salidas derivadas y nuevos conjuntos de datos útiles en el desarrollo de sistemas predictivos de calidad del aire.
